\documentclass[../../../../main.tex]{subfiles}

\begin{document}

\section*{東大理科数学 1995年}

\subsection*{第1問}

\noindent
\textbf{[解]} コーシー・シュワルツより
\[
\frac{\sqrt{2x}}{\sqrt{2}} + 1\sqrt{y} \le \sqrt{\left(\frac{1}{2} + 1\right)(2x+y)} = \sqrt{\frac{3}{2}}\sqrt{2x+y}
\]
等号成立は $\left(\frac{1}{\sqrt{2}}\right) / (\sqrt{2x}) = 1 / (\sqrt{y})$ つまり $x = \frac{1}{4}y$ の時だから, もとめるのは
\[
k = \sqrt{\frac{3}{2}} = \frac{\sqrt{6}}{2}
\]

\newpage

\subsection*{第2問}

\noindent
\textbf{[解]} 対称性から, $y \ge 0$ として良い. $y=0$ の時は等号が成立する. 以下 $y > 0$ とする.
\[
g(x) = \int_0^x (x-t)(1 - \sin t) \, dt
\]
\[
= x \int_0^x (1 - \sin t) \, dt - \int_0^x t(1 - \sin t) \, dt
\]
\[
g'(x) = \int_0^x f(t) \, dt
\]
\[
g''(x) = f(x) \ge 0
\]
より, $g(x)$ は下に凸だから, 凸不等式より
\[
g(x+y) + g(x-y) \ge 2g(x)
\]
である. \hfill $\qed$

\newpage

\subsection*{第3問}

\noindent
\textbf{[解]}

\noindent
(1) $A_n$ のうち, 右端がどのタイルかで場合分けして
\[
A_n = A_{n-2} + (A_{n-1} + A_{n-2}) = 2A_{n-2} + A_{n-1}
\]

\noindent
(2) $A_1 = 1, A_2 = 3$ 及び (1) から
\[
A_n = \frac{1}{3} \left\{ (-1)^n + 2^{n+1} \right\}
\]

\end{document}
